\section{Unknown local directions, nuisance, and the least favorable
entropy
tilt}\label{unknown-local-directions-nuisance-and-the-least-favorable-entropy-tilt}

Part of the \href{DRAFT.md}{manuscript}, continuing
\href{local_prediction.md}{the common-risk construction} and
\href{gaussian_decision.md}{the Gaussian decision reduction}. These are
supplied proofs for a known baseline and uniformly bounded local
perturbations. Sections 1--5 concern a known baseline. Section 6 extends
the analysis to a finite baseline sample in a fixed regular density
neighborhood; arbitrary distribution shifts and joint vanishing-step
limits remain outside scope.

\subsection{1. The model-specific influence
function}\label{the-model-specific-influence-function}

Use the same baseline Q0, fixed x0=0, safe interval {[}-b,b{]}, fixed
positive eta, and entropy update U as in L.1--L.4. Write

\[
U_0(z)=U(0,z),\quad \lambda=E_0\partial_xU(0,Z)>0,
\quad \sigma_g^2=E_0U_0^2>0,
\] \[
s^2=\frac{\eta^2\sigma_g^2}{(1+\eta\lambda)^2-1},
\qquad \psi(z)=\frac{U_0(z)}{\lambda s},
\qquad c^2=E_0\psi^2=\frac{2+\eta\lambda}{\eta\lambda}.
\tag{N.1}
\]

The notation psi is distinct from the logistic variance function a used
in U. It is the influence function of the normalized root displacement
-q(Q)/s at Q0, where s is its fixed baseline value.

\subsubsection{Proposition N.1. Root sensitivity per unit Fisher
information}\label{proposition-n.1.-root-sensitivity-per-unit-fisher-information}

Let h be any bounded, mean-zero score with I\_h=E0 h\^{}2\textgreater0.
Along f\_epsilon=f0(1+epsilon h),

\[
q'_h(0)=-\frac{E_0[U_0h]}{\lambda},
\qquad d_h:=\left.\frac{d}{d\epsilon}\frac{-q(f_\epsilon)}s\right|_0
=E_0[\psi h].
\tag{N.2}
\]

The Fisher information for epsilon in one unlabeled observation is I\_h,
and

\[
\boxed{\frac{d_h^2}{I_h}\le c^2,}
\tag{N.3}
\]

with equality precisely when h is a nonzero scalar multiple of U0 in
L2(Q0).

\textbf{Proof.} Differentiate E\_epsilon U(q\_epsilon,Z)=0. At epsilon=0
the state derivative is lambda and the distribution derivative is E0 U0
h, proving (N.2). Differentiating log f\_epsilon at zero gives score h
and information E0 h\^{}2. Cauchy--Schwarz gives (E0 psi
h)\^{}2\textless=E0 psi\^{}2 E0 h\^{}2. Its equality condition proves
the last assertion. \(\square\)

Thus the original h0=U0/sigma\_g is not merely a convenient tilt: it
maximizes normalized seed displacement for a given local statistical
distinguishability. The term ``least favorable'' here concerns
estimation of this target, after fixing the target window; it is not an
ordering of all nonlinear environmental shifts.

If theta=u/sqrt(n) and the score h is known, its limiting seed is d\_h
u, while the calibration information is rho I\_h. At a fixed seed window
t=d\_h u in {[}-K,K{]}, the resulting Gaussian noise standard deviation
is

\[
r_h=\frac{|d_h|}{\sqrt{\rho I_h}}\le\frac c{\sqrt\rho}.
\tag{N.4}
\]

For d\_h=0 the first-order seed is zero, so this reparameterization is
not applicable; a second-order analysis would be a different local
problem. Comparing different h at a fixed u-window would change the
attainable target window and would not be the same least-favorability
comparison.

\subsection{2. An infinite-dimensional family with an unknown
direction}\label{an-infinite-dimensional-family-with-an-unknown-direction}

Fix finite B\_loc\textgreater0 and K\textgreater0, and define

\[
\mathcal K_{B_{\rm loc},K}
=\{k:E_0k=0,\ \|k\|_\infty\le B_{\rm loc},\ |E_0\psi k|\le K\}.
\tag{N.5}
\]

For each n and each unknown k in this class let

\[
f_{n,k}=f_0(1+k/\sqrt n),\qquad t(k)=E_0\psi k.
\]

The predictor knows f0, U, baseline constants, B\_loc and K. It does not
know k, t(k), the true root, or future SGD samples. The class contains
arbitrary measurable bounded perturbation directions, not just a known
finite list or a fixed parametric span. It is legitimate for the
worst-case k to vary with n.

Keep the same signed density g(z)=kappa tanh(z)f0(z) and joint law

\[
p_{n,k}(z,\pm1)=\tfrac12(f_{n,k}(z)\pm g(z)).
\]

When B\_loc/sqrt(n)\textless=min(1/2,1-kappa), these are legal
normalized densities. Every member shares the complete risk function and
exactly the same global safe interval {[}-b,b{]}. Equal class priors
also remain unchanged. The fixed source information and initialization
remain identical across the enlarged class.

\subsubsection{Lemma N.2. Uniform real-SGD exit law over unknown
directions}\label{lemma-n.2.-uniform-real-sgd-exit-law-over-unknown-directions}

For sufficiently large n, all these fields have one root q\_\{n,k\} on
the same compact exit domain, and uniformly over k,

\[
q_{n,k}=-\frac{E_0[U_0 k]}{\lambda\sqrt n}+O(n^{-1}),
\quad \lambda_{n,k}=\lambda+O(n^{-1/2}),
\quad s_{n,k}=s+O(n^{-1/2}).
\tag{N.6}
\]

Their true right risk-exit probabilities obey

\[
\boxed{\sup_{k\in\mathcal K_{B_{\rm loc},K}}
|\pi_n(Q_{n,k})-\Phi(t(k))|\le C e_n,}
\quad e_n=n^{-1/2}+((1+\log n)/n)^{1/3}.
\tag{N.7}
\]

\textbf{Proof.} For every fixed derivative order j needed by the
stopped-seed proof,

\[
\|v_{n,k}-v_0\|_{C^j}
\le \frac{B_{\rm loc}}{\sqrt n}
\max_{0\le l\le j}\sup_{x,z}|\partial_x^lU(x,z)|.
\]

Thus the entire class lies in one shrinking smooth-field neighborhood.
No derivative of k with respect to z is needed. The baseline derivative
lower bound and endpoint drift signs persist uniformly. Since
v\_\{n,k\}(0)=E0 U0 k/sqrt(n), Taylor expansion at zero and the common
positive derivative give q=O(n\^{}\{-1/2\}) and the root expansion
(N.6). Bounded derivatives and bounded U\^{}2 give the remaining
expansions. The baseline root noise is nonzero, so its uniform lower
bound persists.

For completeness, the probabilistic part can be repeated with exactly
the uniform constants of L.2 rather than applying a fixed-law theorem.
All maps have inverse contraction and C2 conjugacy bounds in the same
smooth neighborhood; their root-to-boundary distances and conjugate exit
magnitudes are bounded away from zero. Bounded U and its derivative give
common batch second/fourth moments and the state-difference noise bound.
For the stopped normalized conjugacy Y and frozen-root seed W, the same
orthogonality and Taylor estimates yield

\[
\sup_{k,j}E|Y_j|^2\le C,\quad
\sup_k E A_{n,k}^{-2\tau}\le C/n,\quad
\sup_k E|Y_\tau-W|^2\le C(1+\log n)/n.
\]

Exit is finite for every member because f\_\{n,k\}\textgreater=f0/2
supplies a common positive mass in each of the fixed forcing intervals.
Frozen seed columns have uniformly bounded third moments and variances
bounded below, so Berry--Esseen and seed anti-concentration hold
uniformly. Choosing the same anti-concentration cutoff as in L.2 bounds
the error of identifying the final exit sign by Ce\_n.~Finally,

\[
\sqrt n H_{n,k}(0)/s_{n,k}=t(k)+O(n^{-1/2})
\]

uniformly by (N.6) and the normalized C2 bound. The global risk-safe
interval is unchanged, so this is actual risk exit without a
component-transfer correction. This proves (N.7). \(\square\)

The proof uses a uniform bounded envelope and a finite number of smooth
U derivatives. It does not require entropy or covering-number control of
the function class of directions, since no uniform empirical process
over k is estimated.

\subsection{3. A direction-free estimator and the minimax
value}\label{a-direction-free-estimator-and-the-minimax-value}

Define the scalar statistic from independent calibration scores

\[
\widehat t=\frac{\sqrt n}{M}\sum_{j=1}^M\psi(\widetilde Z_j).
\tag{N.8}
\]

It is available without knowing k. Under Q\_\{n,k\},

\[
E\widehat t=t(k),\qquad
\operatorname{Var}(\psi(\widetilde Z))
=c^2+n^{-1/2}E_0[\psi^2 k]-t(k)^2/n
=c^2+O(n^{-1/2})
\tag{N.9}
\]

uniformly. Boundedness of psi gives a uniform triangular-array CLT. For
M/n-\textgreater rho in (0,infinity), the distribution of hat t is
uniformly approximated in Kolmogorov distance by N(t(k),c\^{}2/rho).

Let R\^{}K\_\{n,M\} be the minimax expected absolute prediction loss
over (N.5), using the exact pi\_n as the target. Let

\[
V_G(K,r)=\inf_d\sup_{|t|\le K}E|d(t+rZ)-\Phi(t)|.
\]

\subsubsection{Theorem N.3. Unknown-direction local minimax
equivalence}\label{theorem-n.3.-unknown-direction-local-minimax-equivalence}

Assume

\[
B_{\rm loc}\ge K\|\psi\|_\infty/c^2.
\tag{N.10}
\]

Then, for M/n-\textgreater rho in (0,infinity),

\[
\boxed{R^{\mathcal K}_{n,M}\longrightarrow V_G(K,c/\sqrt\rho).}
\tag{N.11}
\]

Uniformly consistent prediction over this enlarged class is possible if
and only if M/n-\textgreater infinity.

\textbf{Proof of the lower bound.} The unknown-direction class contains
the subfamily

\[
k_t=(t/c^2)\psi=(t/c)h_0,\qquad |t|\le K,
\]

by (N.10). Its seed coordinate is exactly t, and it is the original
entropy-gradient tilt with theta=t/(c sqrt(n)). Restricting to this
subfamily and even telling the predictor that this restriction holds can
only reduce the minimax risk. The experiment-level Gaussian reduction
already proved for that subfamily gives liminf
R\^{}K\textgreater=V\_G(K,c/sqrt(rho)). This invokes a full-experiment
lower bound, not a lower bound inferred from a moment CLT.

\textbf{Proof of the upper bound.} Any fixed bounded Lipschitz Gaussian
rule d, applied to hat t, has uniformly convergent risk by (N.9), its
uniform CLT, and (N.7). This suffices to approach V\_G arbitrarily
closely. To see why arbitrary measurable Gaussian rules do not create a
gap, start with an epsilon-optimal rule valued in {[}0,1{]}. Smooth it
by d\_a(x)=E{[}d(x+aZ'){]}, with independent standard-normal Z'. The
smoothed rule is bounded Lipschitz. Jensen gives

\[
E_t|d_a(X)-\Phi(t)|
\le E_t|d(X+aZ')-\Phi(t)|.
\]

The last expectation differs from that under
X\textasciitilde N(t,r\^{}2) by at most
TV(N(0,r\textsuperscript{2+a}2),N(0,r\^{}2)), uniformly in t. This tends
to zero as a-\textgreater0. Thus bounded Lipschitz rules are sufficient
to approximate the Gaussian minimax value. Apply their uniform risk
convergence, then let a and epsilon tend to zero. The upper and lower
bounds agree.

For sufficiency when M/n-\textgreater infinity, use Phi(hat t). Its
error is bounded uniformly by

\[
\frac1{\sqrt{2\pi}}\sqrt{\frac nM}
\sup_k\sqrt{\operatorname{Var}_{n,k}\psi}+Ce_n\to0.
\]

For necessity, a bounded-ratio subsequence contains the original
nondegenerate least-favorable subfamily, to which the two-point lower
bound applies. \(\square\)

This theorem equates a particular decision problem's minimax value. It
does \textbf{not} assert that the entire infinite-dimensional experiment
is equivalent to a scalar Gaussian observation: the full sample may
contain information about orthogonal nuisance directions. The lower
bound comes from a subexperiment and the upper bound from a
direction-free rule.

If M/n-\textgreater0, the same subfamily lower bound and the constant
forecast upper bound also give R\^{}K\_\{n,M\}-\textgreater Phi(K)-1/2.
This extends the earlier zero-information constant to the enlarged
class, without characterizing the entire nuisance experiment.

The fixed finite step rules already computed for the Gaussian problem
also transfer uniformly: their finitely many discontinuities have
vanishing neighborhood probability under the limiting normal densities.
Consequently the existing numerical minimax brackets are relevant to the
enlarged local class without rerunning the optimizer or knowing k.

\subsection{4. What nuisance changes, and what it does
not}\label{what-nuisance-changes-and-what-it-does-not}

\subsubsection{Proposition N.4. Tangent-space
projection}\label{proposition-n.4.-tangent-space-projection}

For a finite collection of bounded centered scores H=(h1,\ldots,hd) with
nonsingular information I=E0 HH\^{}T, let d=E0{[}H psi{]}. The usual
local variance associated with estimating the seed functional in this
score span is

\[
\frac1\rho d^\top I^{-1}d
=\frac1\rho\|\Pi_{\operatorname{span}(H)}\psi\|_2^2
\le\frac{c^2}{\rho}.
\tag{N.12}
\]

\textbf{Proof.} The projection has coefficients I\^{}\{-1\}d, so its
squared L2 norm is d\^{}T I\^{}\{-1\}d.~Projection contracts norm.
Equivalently, an empirical mean of d\^{}T I\^{}\{-1\}H estimates the
derivative functional in this local model with that first-order
variance. \(\square\)

For a closed admissible tangent space T, the corresponding sensitivity
supremum is

\[
\sup_{h\in T,\,\|h\|_2>0}
\frac{\langle\psi,h\rangle^2}{\|h\|_2^2}
=\|\Pi_T\psi\|_2^2.
\tag{N.13}
\]

Equality as an operational bounded density tilt requires that the
maximizing direction be admissible and bounded, or be approximable by
admissible bounded scores. Equation (N.13) alone does not supply the
geometry and uniform exit proof for a new model class.

If an interest score h is accompanied by nuisance scores spanning N, and
those nuisance directions preserve the seed functional to first order,
then N is orthogonal to psi. The efficient interest score is
h\_eff=h-Pi\_N h and

\[
\frac{\langle\psi,h\rangle^2}{\|h_{\rm eff}\|_2^2}
=\frac{\langle\psi,h_{\rm eff}\rangle^2}{\|h_{\rm eff}\|_2^2}
\le c^2.
\tag{N.14}
\]

Thus nuisance can reduce information relative to a restricted
one-direction model, but cannot exceed the full tangent variance c\^{}2
in this known-baseline first-order problem. In particular, for h=h0
proportional to psi and root-preserving nuisance, h\_eff=h0: there is no
additional first-order penalty.

If a proposed nuisance direction has E0 psi k nonzero, it changes the
target and is not a target-preserving nuisance. It must be included in
t(k), as in Theorem N.3, rather than projected out while keeping the
target fixed. Moreover, a bounded multidimensional parameter window can
permit decision improvements beyond a scalar variance comparison; (N.12)
is an information calculation, not a general finite-window minimax
formula.

\textbf{Concrete nuisance example.} At the symmetric Gaussian baseline,
let b(z)=cos(z)-E0 cos(Z). It is bounded, centered and even, whereas psi
and h0 are odd. Thus E0 psi b=E0 h0 b=0. The two-direction perturbation
k\_\{t,v\}=(t/c)h0+v b has seed t, with an unknown bounded v changing
other aspects of the score distribution. It stays in (N.5) whenever
\textbar t\textbar{} H\_*/c+\textbar v\textbar{}
\textbar\textbar b\textbar\textbar\_infinity\textless=B\_loc. This
supplies a legal common-risk nuisance family with no additional
first-order penalty. No new nuisance simulation is claimed in the
present confirmation run.

\subsection{5. Contribution and
boundary}\label{contribution-and-boundary}

The Cauchy--Schwarz and tangent-space projection principles are
classical semiparametric efficiency geometry. The added model-specific
content is the identification psi=U0/(lambda s), its connection to a
uniform true-SGD risk-exit functional, the shared-risk realization of
the bounded perturbation class, and the minimax extension beyond a known
direction.

This result is local around a known Q0. If Q0, its root, lambda or s
must be learned from a separate source sample, that sample's accuracy
and its relation to n enter a new information budget. Arbitrary unknown
baselines, perturbations outside the n\^{}\{-1/2\} neighborhood, growing
B\_loc, online trajectory observation and eta-\textgreater0 are not
covered.

Classical context: van der Laan and Gruber (2016), \emph{One-Step
Targeted Minimum Loss-based Estimation Based on Universal Least
Favorable One-Dimensional Submodels}, discusses scores aligned with
efficient influence curves. This manuscript does not claim that
principle or TMLE methodology as a new contribution.
https://doi.org/10.1515/ijb-2015-0054

\subsection{6. Finite baseline information and conditional
prediction}\label{finite-baseline-information-and-conditional-prediction}

The actual baseline is now unknown within a fixed regular density
neighborhood of a known reference. An independent baseline sample
selects the initialization; the target conditions on that realized
initialization. The complete common risk is retained. Results F.1--F.11
and their proofs follow. Their Gaussian checks do not constitute
actual-SGD or BN/Tent confirmation.

\paragraph{1. An unknown-baseline class with a common
risk}\label{an-unknown-baseline-class-with-a-common-risk}

Use the entropy velocity from the main manuscript:

\[
U(x,z)=-w^2(z-x)a(w(z-x)),\qquad
a(u)=s(u)(1-s(u)),\qquad s(u)=(1+e^{-u})^{-1}.
\]

The confidence slope \(w>0\) and step size \(\eta>0\) are fixed. Write
\(B_U,L_U\) for uniform bounds on \(|U|\) and \(|\partial_xU|\). All
fixed-order state derivatives are uniformly bounded.

Let \(f_\star\) be the centered Gaussian reference density used in the
existing construction. The reference density is known; the actual
baseline is not. Fix \(0<\kappa<1\) and, throughout this chapter, fix

\[
g(z)=\kappa\tanh(z)f_\star(z).
\]

For a sufficiently small fixed \(\delta>0\), let

\[
\mathcal P_\delta
=\left\{P:f_P=f_\star(1+b),\
 E_\star b=0,\ \|b\|_\infty\le\delta\right\}.
\tag{F.1}
\]

The unknown measurable function \(b\) need not be smooth. Smoothness of
the drift comes from \(U\), not from smoothness of \(b\).

For \(P\in\mathcal P_\delta\), define

\[
v_P(x)=E_PU(x,Z),\quad
v_P(q_P)=0,\quad
\lambda_P=v_P'(q_P),\quad
\sigma_P^2=E_PU(q_P,Z)^2,
\]

\[
A_P=1+\eta\lambda_P,\qquad
s_P^2=\frac{\eta^2\sigma_P^2}{A_P^2-1},\qquad
\psi_P(z)=\frac{U(q_P,z)}{\lambda_Ps_P},\qquad
c_P^2=E_P\psi_P^2
      =\frac{2+\eta\lambda_P}{\eta\lambda_P}.
\tag{F.2}
\]

The selected root \(q_P\) is the unique root in a known small interval
specified below. The parameter \(s_P\) is the amplified future-noise
scale, not the logistic function \(s(\cdot)\).

For fixed \(B<\infty\) and \(K\ge0\), let

\[
\mathcal K_{B,K}(P)
=\left\{k:E_Pk=0,\ \|k\|_\infty\le B,\
 |t_P(k)|\le K\right\},\qquad
t_P(k)=E_P[\psi_Pk].
\tag{F.3}
\]

The target marginal is

\[
f_{Q_{n,P,k}}=f_P(1+k/\sqrt n).
\tag{F.4}
\]

For all sufficiently large \(n\), define the labeled target environment
by

\[
p_{n,P,k}(z,\pm1)
=\tfrac12\{f_{Q_{n,P,k}}(z)\pm g(z)\}.
\tag{F.5}
\]

The same construction with \(f_P\) defines a labeled baseline if one is
desired; baseline labels are not observed.

\subparagraph{Lemma F.1. Uniform geometry and legal common
risk}\label{lemma-f.1.-uniform-geometry-and-legal-common-risk}

There exist \(d>0\), \(0<r<d/8\), \(\delta>0\), a fixed admissible step
size, and positive constants independent of \(P,k,n\), such that:

\begin{enumerate}
\def\labelenumi{\arabic{enumi}.}
\tightlist
\item
  Every \(P\in\mathcal P_\delta\) has a unique root
  \(q_P\in[-r/2,r/2]\), and \(v_P'\) has a uniform positive lower bound
  on \([-d,d]\).
\item
  The quantities \(\lambda_P,\sigma_P,s_P,c_P\) are bounded above and
  away from zero. The same root and derivative bounds, with slightly
  enlarged constants, hold for \(Q_{n,P,k}\).
\item
  The densities (F.5) are legal, have equal class priors, and share the
  entire classification-risk function \[
  R(x)=\tfrac12+\int_{-\infty}^{x}g(z)\,dz.
  \tag{F.6}
  \]
\item
  Put \(b_0=d/2\) and \(\varepsilon=R(b_0)-R(0)\). For every
  initialization \(x\in[-r,r]\), the risk-safe set relative to that
  initialization is \[
  \{y:R(y)\le R(x)+\varepsilon\}
  =[-b(x),b(x)],\qquad b_0\le b(x)<3d/4.
  \tag{F.7}
  \] The root is uniformly separated from both boundaries. Choose \[
  0<\eta<\min\{1/L_U,\ d/(8B_U)\}.
  \tag{F.8}
  \] Every one-step overshoot from the safe interval remains in
  \([-d,d]\).
\item
  Uniformly over \(P,k\), \[
  q_{Q_{n,P,k}}
  =q_P-\frac{s_Pt_P(k)}{\sqrt n}+O(n^{-1}),\qquad
  s_{Q_{n,P,k}}=s_P+O(n^{-1/2}).
  \tag{F.9}
  \]
\end{enumerate}

\textbf{Proof.} At \(f_\star\), the root is zero and the root derivative
is strictly positive. Choose \(d\) so the reference drift derivative is
positive throughout \([-d,d]\). For \(j=0,1,2,3\),

\[
\sup_x|v_P^{(j)}(x)-v_\star^{(j)}(x)|
\le\delta\sup_{x,z}|\partial_x^jU(x,z)|.
\]

Choose \(r<d/8\) and then \(\delta\) small enough that the drift is
strictly increasing on \([-d,d]\), its signs at \(\pm r/2\) are
opposite, and its derivative has a uniform positive lower bound. This
gives the unique selected root. Bounded derivatives and the
density-ratio bound give continuity of root noise variance at the
reference law; reduce \(\delta\) to keep it uniformly positive. Formula
(F.2) gives the other bounds.

Choose \(\delta<\min\{1/4,1-\kappa\}\). For large \(n\),

\[
f_Q\ge(1-\delta)(1-B/\sqrt n)f_\star
     \ge |g|.
\]

Thus (F.5) is nonnegative. Since \(\int g=0\), both class priors equal
\(1/2\). The threshold-classification error has derivative \(g\),
proving (F.6). This risk is even and strictly increasing in \(|x|\).
Reduce \(r\), if necessary, until

\[
R(r)-R(0)<R(3d/4)-R(d/2).
\]

If this last step reduces r, reduce delta once more so that the signs at
the new endpoints \(\pm r/2\) still bracket the root; all previous
strict bounds persist on the smaller density class. This fixes the order
of choosing the state and density neighborhoods. This proves (F.7).
Since \(|x_{j+1}-x_j|\le\eta B_U\), (F.8) gives the overshoot statement.

Finally,

\[
v_Q(q_P)=n^{-1/2}E_P[U(q_P,Z)k(Z)]
        =n^{-1/2}\lambda_Ps_Pt_P(k).
\]

The derivative of \(v_Q-v_P\) is uniformly \(O(n^{-1/2})\). A uniform
Taylor expansion of the root equation gives the first part of (F.9); the
root, variance and slope expansions give the second. \(\square\)

\textbf{Information distinction.} The known risk identifies the fixed
signed density \(g\), and may identify the known reference \(f_\star\).
It does not identify the unknown multiplier \(b\) in (F.1). In
particular, \(g\) must not be redefined as \(\kappa\tanh(z)f_P(z)\) when
\(P\) changes.

\paragraph{2. Observations, initialization and conditional
target}\label{observations-initialization-and-conditional-target}

The predictor receives independent samples

\[
B_L=(B_1,\ldots,B_L)\sim P^{\otimes L},\qquad
C_M=(C_1,\ldots,C_M)\sim Q_{n,P,k}^{\otimes M}.
\tag{F.10}
\]

Both samples are independent of all future SGD batches. The predictor
knows the declared class, \(U,\eta,w,R,B,K\), and the initialization
rule, but not \(P,k,q_P,q_Q\), or their exact dynamical constants.

Define

\[
\widehat v_L(x)=L^{-1}\sum_{i=1}^{L}U(x,B_i).
\]

If \(\widehat v_L\) has a strictly positive derivative throughout
\([-r,r]\), bounded below by half a known uniform drift-derivative lower
bound, and has opposite signs at the endpoints, let \(x_L\) be its
unique root in that interval. Otherwise set \(x_L=0\). This is a
measurable, specified fallback rule; failed root estimation is not
discarded.

Future adaptation is

\[
x_0=x_L,\qquad
x_{j+1}=x_j+\frac{\eta}{n}\sum_{i=1}^{n}U(x_j,Z_{j,i}),
\qquad Z_{j,i}\stackrel{\mathrm{iid}}{\sim}Q_{n,P,k}.
\tag{F.11}
\]

Let

\[
T_\varepsilon
=\inf\{j\ge0:R(x_j)>R(x_L)+\varepsilon\}.
\]

By Lemma F.1 this is exactly exit from \([-b(x_L),b(x_L)]\), including
overshoots. Define the conditional forecast target

\[
\pi_n(P,k;B_L)
=\Pr_{Q_{n,P,k}}\{x_{T_\varepsilon}>b(x_L)\mid B_L\}.
\tag{F.12}
\]

The expectation in (F.12) is only over the future trajectory.
Conditional on \(B_L,C_M\), the future probability is the same number
because the calibration sample is independent of the future.

The minimax MAE is

\[
\mathfrak R_{n,L,M}
=\inf_{\widehat\pi}
 \sup_{\substack{P\in\mathcal P_\delta\\k\in\mathcal K_{B,K}(P)}}
 E_{P,Q}\left|\widehat\pi(B_L,C_M)-\pi_n(P,k;B_L)\right|.
\tag{F.13}
\]

Actions can be restricted to \([0,1]\). Independent randomization cannot
improve this convex loss.

\paragraph{3. Root estimation and actual-SGD
exit}\label{root-estimation-and-actual-sgd-exit}

\subparagraph{Lemma F.2. Uniform initialization expansion and scale
estimation}\label{lemma-f.2.-uniform-initialization-expansion-and-scale-estimation}

Uniformly over \(P\in\mathcal P_\delta\), the fallback probability is at
most \(C e^{-cL}\), and

\[
x_L-q_P
=-\frac1{\lambda_PL}\sum_{i=1}^{L}U(q_P,B_i)+R_{L,P},
\qquad
\sup_P E_P|R_{L,P}|\le C/L,
\tag{F.14}
\]

\[
\sup_P E_P|x_L-q_P|^2\le C/L,\qquad
\sup_P E_P|x_L-q_P|^4\le C/L^2.
\tag{F.15}
\]

Use the baseline sample to estimate

\[
\widehat\lambda=\widehat v_L'(x_L),\qquad
\widehat\sigma^2=L^{-1}\sum_i
\{U(x_L,B_i)-\widehat v_L(x_L)\}^2,
\]

and truncate these estimates to fixed positive compact intervals
containing their true ranges in their interiors. Define
\(\widehat s,\widehat c\) through (F.2). Then these estimates are
uniformly consistent, and their squared errors are \(O(L^{-1})\).

\textbf{Proof.} Bounded \(U\) gives exponential concentration at both
bracket endpoints. Bounded \(\partial_x^2U\) permits a fixed finite grid
on \([-r,r]\) for testing a fixed positive derivative margin; Hoeffding
bounds on that grid give the uniform exponential bound for derivative
failure.

On the good event, the empirical inverse derivative is bounded. The mean
value theorem gives

\[
|x_L-q_P|\le C|\widehat v_L(q_P)|.
\]

The centered bounded sample mean has second moment \(O(L^{-1})\) and
fourth moment \(O(L^{-2})\). The fallback contributes exponentially
small terms because \(x_L,q_P\) are bounded. Taylor expansion gives

\[
0=\widehat v_L(q_P)+\lambda_P(x_L-q_P)
 +\{\widehat v_L'(q_P)-\lambda_P\}(x_L-q_P)
 +O(|x_L-q_P|^2).
\]

Cauchy--Schwarz and the preceding moment bounds give (F.14). For the
squared remainder, write the derivative fluctuation at \(q_P\) as
\(D_L\). Its fourth moment is \(O(L^{-2})\); hence
\(E[D_L^2(x_L-q_P)^2]=O(L^{-2})\) by Cauchy--Schwarz and (F.15). The
quadratic Taylor term has the same squared bound. The bad event has
exponential probability and all relevant quantities are bounded, so
\(E|R_{L,P}|^2=O(L^{-2})\) uniformly as well. For scale estimates, first
evaluate the empirical moments at \(q_P\), then use bounded derivatives
to replace \(q_P\) by \(x_L\). Smoothness of (F.2) on the fixed positive
compact ranges completes the proof. \(\square\)

\subparagraph{Lemma F.3. Uniform exit law with a varying initial
point}\label{lemma-f.3.-uniform-exit-law-with-a-varying-initial-point}

For each fixed finite \(H\), uniformly over \(P,k\) and deterministic
starting points \(x\in[-r,r]\) satisfying

\[
\sqrt n\,|x-q_Q|\le H,
\]

the actual fresh-batch dynamics satisfy

\[
\left|
\Pr_Q\{x_{T_\varepsilon}>b(x)\mid x_0=x\}
-\Phi\!\left(\frac{\sqrt n(x-q_Q)}{s_Q}\right)
\right|\le C_H e_n,
\tag{F.16}
\]

where

\[
e_n=n^{-1/2}+\left(\frac{1+\log n}{n}\right)^{1/3}.
\]

The constants are uniform over the full bounded density class; no
differentiability in the distribution parameter is needed.

\textbf{Proof.} Here \(Q=Q_{n,P,k}\). All laws have a common density
lower bound relative to \(f_\star\). A fixed bounded score interval
strictly to the left of \([-d,d]\) gives a uniformly positive entropy
velocity throughout the state interval. A fixed number of consecutive
batches whose scores all fall in that interval forces exit. The block
probability is positive for each \(n\), uniformly in the state, so
\(T_\varepsilon<\infty\) almost surely. This argument does not supply a
typical exit-time estimate.

Put \(G_Q(y)=y+\eta v_Q(y)\), \(A_Q=1+\eta v_Q'(q_Q)\). The inverse map
contracts on \([-d,d]\). The inverse-iteration construction gives a
normalized increasing \(C^2\) conjugacy \(H_Q\) such that

\[
H_Q(q_Q)=0,\quad H_Q'(q_Q)=1,\quad
H_Q(G_Q(y))=A_QH_Q(y)
\]

where both arguments lie in the interval. Uniform bounds are

\[
\|H_Q\|_{C^2}\le C,\quad \inf H_Q'\ge m>0,\quad
1<A_{\min}\le A_Q\le A_{\max}.
\tag{F.17}
\]

For completeness, center the inverse map at \(q_Q\), write it as
\(A_Q^{-1}y+O(y^2)\), and construct \(H_Q\) as the limit of
\(A_Q^j(G_Q^{-j}(y)-q_Q)\). Reduce the fixed density neighborhood and
the common state neighborhood so that the centered inverse derivatives
are bounded by a common \(\beta\) with \(A_{\max}\beta^2<1\). This is
possible because at the reference root
\(A_\star(A_\star^{-1})^2=A_\star^{-1}<1\), and roots and derivatives
vary uniformly continuously under the density envelope. The differences
of successive normalized inverse iterates and their first two
derivatives are then bounded by a summable geometric sequence. The
second-derivative difference is controlled by the uniform third
derivative of the inverse and the contraction of the iterates. Uniform
bounded third derivatives of the drift supply the second-derivative
control. The derivative product has positive factors
\(1+O(|G_Q^{-j}(y)-q_Q|)\); the sum of deviations is bounded, so its
limit is bounded away from zero. A fixed number of inverse iterations
covers the full compact interval. This proves (F.17) uniformly. These
are exactly the ingredients used in the original actual-SGD argument,
now uniform in \(P,k\) and in the initial point.

Let \(\tau=T_\varepsilon\),
\(\xi_j(y)=n^{-1}\sum_iU(y,Z_{j,i})-v_Q(y)\), and \[
Y_t=\sqrt n A_Q^{-(t\wedge\tau)}H_Q(x_{t\wedge\tau}),\quad
\mu=\sqrt nH_Q(x).
\]

Taylor expansion before exit yields a martingale increment plus a
remainder bounded by \(C\xi_j(x_j)^2\). Conditional batch second and
fourth moments are bounded by \(C/n\) and \(C/n^2\). Since
\(|\mu|\le C_H\), the stopped recursion gives

\[
\sup_t E|Y_t|^2\le C_H.
\]

The root-to-boundary distance and \(\inf H_Q'>0\) imply \[
EA_Q^{-2\tau}\le C_H/n,\qquad
E[1_{\{j<\tau\}}|x_j-q_Q|^2]
\le C_H\min(A_Q^{2j}/n,1).
\tag{F.18}
\]

Introduce the frozen-root seed on the same infinite sample array, \[
W=\mu+\eta\sum_{j\ge0}A_Q^{-j-1}\sqrt n\,\xi_j(q_Q).
\]

The conditional variance of \[
H_Q'(G_Q(x_j))\xi_j(x_j)-\xi_j(q_Q)
\] is at most \(C|x_j-q_Q|^2/n\). The weighted sum and (F.18) therefore
contribute \(C_H(1+\log n)/n\) to the squared approximation error.
Taylor remainders contribute \(C/n\). The frozen tail after \(\tau\),
conditional on the stopping history, has variance \(s_Q^2A_Q^{-2\tau}\).
Consequently \[
E|Y_\tau-W|^2\le C_H(1+\log n)/n.
\tag{F.19}
\]

Grouping \(W-\mu\) by sample column gives a normalized sum of \(n\) iid
centered variables with uniform bounded third moments and variance
\(s_Q^2\) bounded away from zero. Berry--Esseen gives a uniform
\(C/\sqrt n\) Gaussian approximation and \[
\Pr(|W|\le u)\le Cu+C/\sqrt n.
\]

The sign of \(Y_\tau\) is exactly the exit side. Apply (F.19), Chebyshev
and this anti-concentration bound with \(u=((1+\log n)/n)^{1/3}\).
Finally, \[
\sqrt nH_Q(x)=\sqrt n(x-q_Q)+O_H(n^{-1/2}).
\] This proves (F.16). \(\square\)

\subparagraph{Theorem F.4. Conditional exit and the two-sample
interface}\label{theorem-f.4.-conditional-exit-and-the-two-sample-interface}

Assume \(L/n\to\ell\in(0,\infty]\). Define

\[
S_n=\frac{\sqrt n(x_L-q_Q)}{s_P},\qquad
\widehat D_n=
\frac{\sqrt n}{\widehat\lambda\,\widehat s}\,
\frac1M\sum_{j=1}^{M}U(x_L,C_j).
\tag{F.20}
\]

Then

\[
\sup_{P,k}E_B
\left|\pi_n(P,k;B_L)-\Phi(S_n)\right|\longrightarrow0.
\tag{F.21}
\]

If also \(M/n\to\rho\in(0,\infty)\), the joint law of
\((S_n,\widehat D_n)\) converges, uniformly in bounded-Lipschitz
distance over \(P,k\), to

\[
S=t+a_PZ_L,\qquad D=S+b_PZ_C,\qquad
a_P=\frac{c_P}{\sqrt\ell},\quad
b_P=\frac{c_P}{\sqrt\rho},
\tag{F.22}
\]

where \(t=t_P(k)\), \(Z_L,Z_C\) are independent standard normals, and
\(a_P=0\) when \(\ell=\infty\). If \(M/n\to\infty\), then

\[
\widehat D_n-S_n\longrightarrow0
\]

uniformly in probability, and \(\Phi(\widehat D_n)\) is uniformly
consistent for the conditional target (F.12).

\textbf{Proof.} Lemmas F.1--F.2 give

\[
S_n
=t_P(k)-\frac{\sqrt n}{L}\sum_{i=1}^{L}\psi_P(B_i)+o_P(1),
\tag{F.23}
\]

uniformly. In particular, \(\sup_{P,k}E S_n^2\le C\). On \(|S_n|\le H\),
Lemma F.3 applies conditionally on \(B_L\), since future batches are
independent. On the complement use the bound one. The replacement of
\(s_Q\) by \(s_P\) costs \(o(1)\). First send \(n\) to infinity for
fixed \(H\), then send \(H\) to infinity using the uniform second
moment. This proves (F.21). It is an average bound on errors in
conditional probabilities; it does not replace the conditional target by
its average.

For the calibration statistic, Taylor expansion at \(q_P\) gives, when
\(M/n\) stays bounded away from zero,

\[
\widehat D_n
=t_P(k)-\frac{\sqrt n}{L}\sum_i\psi_P(B_i)
 +\frac{\sqrt n}{M}\sum_j
 \{\psi_P(C_j)-E_Q\psi_P(C_j)\}
 +o_P(1).
\tag{F.24}
\]

To verify the remainder, \(x_L-q_P=O_P(L^{-1/2})\),
\(E_QU_x(q_P)=\lambda_P+O(n^{-1/2})\), and the calibration derivative
fluctuation is \(O_P(M^{-1/2})\). The Taylor quadratic term is
\(O_P(\sqrt n/L)=o_P(1)\). Replacing the denominator by the consistent
estimated denominator is valid because the unstudentized statistic is
tight. All statements are uniform. The two leading centered sums in
(F.24) are independent; their variances are \(c_P^2n/L\) and
\(c_P^2n/M+o(1)\). Uniform boundedness of \(\psi_P\), a uniform positive
variance, and Berry--Esseen give (F.22). If \(L/n\to\infty\), the first
centered sum vanishes.

For \(M/n\to\infty\), conditional calibration variance at \(x_L\) is at
most \(C/M\). The same Taylor bounds give \(\widehat D_n-S_n=o_P(1)\).
Boundedness and the Lipschitz property of \(\Phi\), together with
(F.21), prove consistency. \(\square\)

\textbf{Cancellation.} The realized baseline-root error occurs in both
\(S_n\) and \(\widehat D_n\). Thus the leading error of the direct
predictor relative to its \emph{conditional} target has calibration
variance \(c_P^2n/M\), not \(c_P^2n(1/L+1/M)\). The latter is the
variance of a root-difference statistic around the deterministic shift
\(t\), which is a different estimand.

\paragraph{4. Exact limiting risk of the raw feasible
rule}\label{exact-limiting-risk-of-the-raw-feasible-rule}

Define \[
\mathcal D(u,v)
=E|\Phi(u+vZ)-\Phi(u)|
=\int_0^\infty\{\Phi(u+vz)-\Phi(u-vz)\}\phi(z)\,dz.
\tag{F.25}
\]

\subparagraph{Corollary F.5. Raw-rule risk and nonmonotone baseline
effects}\label{corollary-f.5.-raw-rule-risk-and-nonmonotone-baseline-effects}

Under the finite-ratio assumptions of Theorem F.4,

\[
\sup_{P,k}\left|
E|\Phi(\widehat D_n)-\pi_n(P,k;B_L)|
-
\mathcal D\!\left(
\frac{t_P(k)}{\sqrt{1+a_P^2}},
\frac{b_P}{\sqrt{1+a_P^2}}
\right)
\right|\longrightarrow0.
\tag{F.26}
\]

For each fixed \(P\), the maximum limiting risk over the declared
direction class is

\[
\frac1\pi\arctan\!\left(
\frac{b_P}{\sqrt{1+a_P^2}}\right),
\tag{F.27}
\]

attained at \(t=0\), which the class contains.

\textbf{Proof.} Apply Theorem F.4 to the bounded Lipschitz loss
\(|\Phi(D)-\Phi(S)|\). Conditional on \(S\), its expectation is
\(\mathcal D(S,b_P)\). Integrating over \(S=t+a_PZ_L\) uses \[
E\Phi(t+a_PZ_L+y)
=\Phi\!\left(\frac{t+y}{\sqrt{1+a_P^2}}\right),
\] which proves (F.26). The function \(\mathcal D(u,v)\) is even and
decreasing in \(|u|\); pairing \(z\) and \(-z\) and differentiating
proves this directly. At zero, \(\mathcal D(0,v)=\arctan(v)/\pi\),
obtained by differentiating the integral in \(v\) and integrating from
zero. \(\square\)

For fixed calibration noise \(b_P\), increasing initialization noise
\(a_P\) decreases this average raw-rule risk. This is not evidence that
a poorer initialization is beneficial: the conditional probability is
more often near zero or one. The data-generating target changes with
\(L\).

Also, \[
E_B\pi_n(P,k;B_L)
-\Phi\!\left(\frac{t_P(k)}{\sqrt{1+a_P^2}}\right)\longrightarrow0
\tag{F.28}
\] uniformly. This marginal probability is not the target in (F.13).

\paragraph{5. The Gaussian problem with a random
target}\label{the-gaussian-problem-with-a-random-target}

For \(a,b>0\), let

\[
S=t+aZ_L,\qquad D=S+bZ_C,
\qquad |t|\le K,
\]

with independent standard normals. Define

\[
V_{\mathrm{FB}}(K;a,b)
=\inf_{d:\mathbb R\to[0,1]}
 \sup_{|t|\le K}E_t|d(D)-\Phi(S)|.
\tag{F.29}
\]

Both the observation and the target contain \(Z_L\). The parameter
restriction is on \(t\), not on the realized \(S\).

The corresponding experiment with an unknown baseline position is

\[
X=\beta+aZ_L,\qquad Y=\beta-t+bZ_C,\qquad
\text{target}=\Phi(X-\beta+t).
\tag{F.30}
\]

Let \(V_H(K;a,b)\) be its minimax risk when \(|\beta|\le H\), and let
\(V_\infty\) allow every real \(\beta\).

\subparagraph{Proposition F.6. Removing the unrestricted baseline
position}\label{proposition-f.6.-removing-the-unrestricted-baseline-position}

For \(a,b>0\),

\[
V_\infty(K;a,b)=V_{\mathrm{FB}}(K;a,b),
\tag{F.31}
\]

and, for \(H>0\),

\[
0\le V_{\mathrm{FB}}(K;a,b)-V_H(K;a,b)
\le E\min\!\left(\frac{|aZ|}{2H},1\right)
\le\frac{a}{H\sqrt{2\pi}}.
\tag{F.32}
\]

Thus discarding the absolute coordinate is justified in the
unrestricted-location minimax problem, or after taking \(H\to\infty\).
It need not be optimal for a fixed bounded location window.

\textbf{Proof.} For any rule \(d\), use \(\delta(X,Y)=d(X-Y)\). With
\(S=X-\beta+t\), the pair \((S,X-Y)\) has the law in (F.29), up to the
irrelevant sign of \(Z_C\). Hence \(V_H,V_\infty\le V_{\mathrm{FB}}\).

For the reverse bound, let \(\mu\) be any prior on \([-K,K]\),
independently assign \(\beta\) the uniform prior on \([-H,H]\), and
write \(D=X-Y\), \(\epsilon_L=S-t\). Conditional on \((t,S,D)\), the
observed \(X\) is uniform on \([-H+\epsilon_L,H+\epsilon_L]\).

Compare this joint distribution to one with an auxiliary observation
\(X'\), uniform on \([-H,H]\) and independent of \((t,S,D)\). The total
variation distance is at most \[
E\min(|\epsilon_L|/(2H),1).
\] In the auxiliary problem \(X'\) cannot reduce Bayes risk: averaging a
randomized action over \(X'\) improves or preserves absolute loss by
Jensen. Therefore the original Bayes risk is at least the Bayes risk for
observing \(D\) alone, minus this TV bound.

Taking the supremum over priors \(\mu\) gives (F.32). The minimax
equality between rules and priors is justified below. Letting
\(H\to\infty\), and using the upper bound, proves (F.31). \(\square\)

\subparagraph{Proposition F.7. Bayes characterization, exact subcase and
checks}\label{proposition-f.7.-bayes-characterization-exact-subcase-and-checks}

The following properties hold.

\textbf{(a) Minimax/Bayes characterization.} Let \(\phi_a\) denote the
\(N(0,a^2)\) density. Then

\[
V_{\mathrm{FB}}(K;a,b)
=\sup_{\mu\in\mathcal P([-K,K])}
\int_{\mathbb R}\inf_{p\in[0,1]}
\int_{-K}^{K}\int_{\mathbb R}
|p-\Phi(s)|\phi_a(s-t)\phi_b(d-s)\,ds\,\mu(dt)\,dd.
\tag{F.33}
\]

A least-favorable prior exists. For a given prior, a Bayes rule is the
posterior median of \(\Phi(S)\), equivalently \(\Phi\) of the posterior
median of \(S\).

Put \[
\alpha=\frac{a^2}{a^2+b^2},\qquad
v^2=\frac{a^2b^2}{a^2+b^2}.
\] Conditional on \(D=d,t\), \[
S\mid D=d,t\sim
N\{\alpha d+(1-\alpha)t,\ v^2\},
\] and the posterior mixing weights for \(t\) are proportional to
\(\phi_{\sqrt{a^2+b^2}}(d-t)\mu(dt)\). These formulas specify the
posterior median by a one-dimensional monotone equation.

\textbf{(b) No root displacement.} For \(K=0\), \[
d_0^\star(d)=\Phi\!\left(\frac{a^2}{a^2+b^2}d\right)
\tag{F.34}
\] is minimax, and \[
V_{\mathrm{FB}}(0;a,b)
=\frac1\pi\arctan
\left(\frac{ab}{\sqrt{a^2+b^2+a^4}}\right).
\tag{F.35}
\]

In particular, \(V_{\mathrm{FB}}(0;1,1)=1/6\), although the marginal
probability \(E\Phi(S)\) is exactly \(1/2\).

\textbf{(c) General-window bounds.} For \(K\ge0\), \[
\frac1\pi\arctan
\left(\frac{ab}{\sqrt{a^2+b^2+a^4}}\right)
\le V_{\mathrm{FB}}(K;a,b)
\le \frac1\pi\arctan\left(\frac{b}{\sqrt{1+a^2}}\right).
\tag{F.36}
\] The upper bound is the raw rule \(d(d)=\Phi(d)\). The lower bound
restricts to \(t=0\).

\textbf{(d) Known-baseline recovery and perfect calibration.} Defining
the degenerate cases directly, \[
V_{\mathrm{FB}}(K;0,b)=V_G(K,b),\qquad
V_{\mathrm{FB}}(K;a,0)=0.
\tag{F.37}
\] Moreover, \(V_{\mathrm{FB}}(K;a,b)\to V_G(K,b)\) as \(a\downarrow0\)
for fixed \(b>0\), and \(V_{\mathrm{FB}}(K;a,b)\to0\) as
\(b\downarrow0\).

\textbf{(e) Continuity and finite-grid control.} On compact subsets of
\(a,b>0\), \(V_{\mathrm{FB}}\) is continuous in \((K,a,b)\). For every
rule \(d\), \[
|R_t(d)-R_u(d)|
\le 2\Phi\!\left(\frac{|t-u|}{2a}\right)-1
\le \frac{|t-u|}{a\sqrt{2\pi}}.
\tag{F.38}
\] Consequently a grid with maximum distance \(h\) to \([-K,K]\) gives
\[
\sup_{|t|\le K}R_t(d)
\le\max_{t\in\mathrm{grid}}R_t(d)
     +h/(a\sqrt{2\pi}).
\tag{F.39}
\] This is a continuous-window risk bound; numerical quadrature and tail
errors must still be controlled separately.

\textbf{Proof.} For (a), rule risks are convex in the action function.
Use a common Gaussian reference measure on the observation line. The set
of \([0,1]\)-valued rules is weak-* compact in \(L^\infty\), and
integral convexity makes each risk lower semicontinuous in this
topology. Risks are continuous in \(t\), uniformly over rules, by the TV
bound (F.38). Convex minimax therefore interchanges the rule and prior
optimizations. The prior space is weakly compact, and Bayes risk is an
upper semicontinuous infimum of continuous affine functions of the
prior, giving a least-favorable prior. Pointwise minimization under
absolute loss gives a posterior median. For \(a,b>0\), the posterior
density of \(S\) is strictly positive, so its median is unique.

For (b), \(t=0\) is known in (F.29); Gaussian conditioning gives the
posterior median \(\alpha D\). Write \[
S=m+vZ,\qquad m=\alpha D,\qquad
\operatorname{Var}(m)=d_m^2=\frac{a^4}{a^2+b^2},
\] where \(m\) and \(Z\) are independent centered normals. The risk is
\[
E|\Phi(m+vZ)-\Phi(m)|
=\frac1\pi\arctan\!\left(\frac{v}{\sqrt{1+d_m^2}}\right)
\] by the calculation in Corollary F.5. Simplifying gives (F.35). At
\(a=b=1\), the arctangent argument is \(1/\sqrt3\), giving \(1/6\).

Part (c) follows from restriction to \(t=0\) and Corollary F.5.

For (d), when \(a=0\), \(S=t\) is deterministic and the experiment is
exactly the original Gaussian problem. When \(b=0\), \(S=D\) is
observed. To prove recovery as \(a\downarrow0\), replace the target
\(\Phi(S)\) by \(\Phi(t)\), at a risk cost at most \[
E|\Phi(t+aZ)-\Phi(t)|\le a/\pi.
\] After this replacement, \(D\mid t\sim N(t,a^2+b^2)\), so the minimax
value is \(V_G(K,\sqrt{a^2+b^2})\). Gaussian TV continuity gives
convergence to \(V_G(K,b)\). The raw upper bound proves the
\(b\downarrow0\) statement.

For (e), changing \(t\) changes only the mean in \(S\sim N(t,a^2)\); the
conditional law \(D\mid S\) and the bounded payoff are unchanged. TV
contraction gives (F.38). Gaussian covariance continuity gives uniform
TV continuity in \((a,b)\) away from zero. These bounds pass through the
minimax optimizations and give continuity, including variation of the
compact endpoint \(K\). \(\square\)

\textbf{Why the old window projection fails.} For \(a>0\), \(S\) has
positive probability outside \([-K,K]\). A rule whose outputs stay in
\([\Phi(-K),\Phi(K)]\) retains truncation error even as
\(b\downarrow0\). For instance, at \(t=0,a=1,K=1\), the best rule
subject to that output range has limiting risk \[
E|\Phi(\operatorname{clip}(S,-1,1))-\Phi(S)|
=(1-\Phi(1))^2\approx0.02517149.
\tag{F.40}
\] Here \(\Phi(S)\) is uniform on \([0,1]\), so the identity follows by
integrating the two clipped tails. This is a Gaussian calculation, not
an actual-SGD experiment.

\paragraph{6. A marked two-sample experiment: keeping the
initialization}\label{a-marked-two-sample-experiment-keeping-the-initialization}

Fix an interior baseline \(P_0\in\mathcal P_\delta\), meaning its
density multiplier has norm strictly less than \(\delta\). Write
\(q_0,s_0,c_0,\psi_0\) for its constants.

For finite \(H\) and a seed window \(K'<K\), define, relative to
\(P_0\),

\[
f_{P_{n,\beta}}
=f_{P_0}\left(1-\frac{\beta\psi_0}{c_0^2\sqrt n}\right),
\] \[
f_{Q_{n,\beta,t}}
=f_{P_0}\left(1-\frac{(\beta-t)\psi_0}{c_0^2\sqrt n}\right),
\quad |\beta|\le H,\quad |t|\le K'.
\tag{F.41}
\]

Both integrate to one. The baseline moves its root to first order by
\(s_0\beta/\sqrt n\); the target root is moved by
\(s_0(\beta-t)/\sqrt n\).

Relative to \(P_{n,\beta}\), the target perturbation is

\[
k_{n,\beta,t}(z)
=
\frac{t\psi_0(z)/c_0^2}
 {1-\beta\psi_0(z)/(c_0^2\sqrt n)}.
\tag{F.42}
\]

Its baseline mean is exactly zero. Its seed parameter is
\(t+O(n^{-1/2})\), uniformly on the compact parameter rectangle. Thus it
belongs to the declared class for large \(n\), provided the envelope has
slack for \(K'\). For \(K=0\), take \(t=0\), \(k=0\), without needing an
interior seed window.

\subparagraph{Lemma F.8. Compact marked-experiment minimax
convergence}\label{lemma-f.8.-compact-marked-experiment-minimax-convergence}

Assume \(L/n\to\ell\in(0,\infty)\) and \(M/n\to\rho\in(0,\infty)\). On
the submodel (F.41), the full two-sample prediction problem converges in
minimax value to (F.30) restricted to \(|\beta|\le H,|t|\le K'\), with
\[
a=c_0/\sqrt\ell,\qquad b=c_0/\sqrt\rho.
\tag{F.43}
\]

The transfer preserves the observed initialization in the loss. A CLT
for \(X-Y\) alone is insufficient for this conclusion.

\textbf{Proof.} Put \(h_0=\psi_0/c_0\), so \(E_0h_0=0\), \(E_0h_0^2=1\),
and \(h_0\) is bounded. At the product baseline
\(P_0^{\otimes L}\otimes P_0^{\otimes M}\), define independent
standardized score sums \[
S_L=L^{-1/2}\sum_i h_0(B_i),\qquad
S_M=M^{-1/2}\sum_j h_0(C_j).
\]

Their distributions are atomless: \(U(q_0,\cdot)\) is a nonconstant
analytic function and \(P_0\) has a density. Let \(F_L,F_M\) be their
known baseline CDFs in this fixed submodel, and set \[
T_L=\Phi^{-1}(F_L(S_L)),\qquad
T_M=\Phi^{-1}(F_M(S_M)).
\]

Under the product baseline, these are exactly independent standard
normals. Berry--Esseen, tightness and uniform integrability give
\(S_L-T_L\to0\), \(S_M-T_M\to0\) in probability and in every fixed
\(L^p\) needed here. Exponential moments of the score sums are bounded
uniformly at each fixed exponent, by bounded-score concentration.

The log likelihood on the parameter rectangle has expansion \[
\log\Lambda_{n,\beta,t}
=-\frac{\beta\sqrt{L/n}}{c_0}S_L
 -\frac{(\beta-t)\sqrt{M/n}}{c_0}S_M
 -\frac{\beta^2L/n+(\beta-t)^2M/n}{2c_0^2}
 +o_{P_0}(1),
\tag{F.44}
\] uniformly in \((\beta,t)\). The deterministic third-order Taylor
remainder is \(O((L+M)n^{-3/2})=o(1)\); the empirical quadratic terms
converge uniformly since the parameter rectangle is compact.

Replace the score sums by \(T_L,T_M\) in the exponential of (F.44), and
call the resulting exact Gaussian likelihood \(g_{n,\beta,t}\). Its
baseline expectation is one. Uniform exponential moments and the uniform
log-likelihood remainder imply \[
E_0\sup_{\beta,t}
|\Lambda_{n,\beta,t}-g_{n,\beta,t}(T_L,T_M)|
\longrightarrow0.
\tag{F.45}
\]

One direct uniform-integrability justification is to bound the supremum
of either likelihood by a constant times \(\exp\{C(|S_L|+|S_M|)\}\) or
\(\exp\{C(|T_L|+|T_M|)\}\), whose moments of any fixed order greater
than one are uniformly bounded. Uniform convergence in probability then
yields (F.45).

The auxiliary full-sample measures with densities
\(g_{n,\beta,t}(T_L,T_M)\) relative to the product baseline are
therefore uniformly close in TV to the actual measures. Under the
auxiliary experiment, \[
X_n^G=-c_0\sqrt{n/L}\,T_L,\qquad
Y_n^G=-c_0\sqrt{n/M}\,T_M
\] have exactly the independent Gaussian laws \[
X_n^G\sim N(\beta,c_0^2n/L),\qquad
Y_n^G\sim N(\beta-t,c_0^2n/M).
\tag{F.46}
\] The baseline conditional distribution of the full samples given
\((T_L,T_M)\) reconstructs the auxiliary full-sample experiment and is
parameter independent.

It remains to preserve the loss, since its target is sample dependent.
Let \[
X_n^I=\frac{\sqrt n(x_L-q_0)}{s_0}.
\] Lemma F.2 gives \[
X_n^I=-c_0\sqrt{n/L}\,S_L+o_{P_0}(1)
      =X_n^G+o_{P_0}(1).
\tag{F.47}
\] Uniform likelihood moment bounds transfer the bounded effect of
(F.47) to every parameter on the compact rectangle. Root expansions and
Theorem F.4 give \[
\sup_{\beta,t}E_{\beta,t}
\left|\pi_n(P_{n,\beta},k_{n,\beta,t};B_L)
 -\Phi(X_n^G-\beta+t)\right|\longrightarrow0.
\tag{F.48}
\]

After replacing the target by the right side of (F.48), conditioning any
full-sample action on \((T_L,T_M)\) under the auxiliary experiment
cannot increase risk, by convexity. Its conditional law is parameter
independent. Conversely, every Gaussian rule can be applied to the
observable \((X_n^G,Y_n^G)\) in the original experiment. More
explicitly, let \(\epsilon_n\) be the uniform TV error in (F.45), and
let \(r_n\) be the supremum of the expected target-replacement error in
(F.48). For every {[}0,1{]}-valued full-sample rule, replacing its
target and its law changes each risk by at most \(r_n+\epsilon_n\).
Under the auxiliary law the conditional distribution of the original
samples given the Gaussian pair is the baseline conditional kernel; its
conditional mean action therefore defines one parameter-independent
Gaussian rule. Jensen lowers risk because the replacement target is a
function of that Gaussian pair and the parameter. In the other
direction, composing any Gaussian rule with the Gaussian pair gives a
full-sample rule, with the same \(r_n+\epsilon_n\) bound. Thus the
absolute difference of the two minimax values tends to zero. This
argument transfers the observation together with the initialization mark
in the loss, not only the distribution of the difference statistic.
Equations (F.45) and (F.48) transfer both minimax inequalities. Finally,
covariance continuity replaces the finite ratios in (F.46) by their
positive limits. This proves the asserted convergence. \(\square\)

\paragraph{7. Minimax prediction over the unknown-baseline
class}\label{minimax-prediction-over-the-unknown-baseline-class}

For the lower bound, assume the direction envelope contains the complete
aligned subfamily:

\[
B\ge K\sup_{P\in\mathcal P_\delta}
\frac{\|\psi_P\|_\infty}{c_P^2}.
\tag{F.49}
\]

The right side is finite by Lemma F.1. For \(K=0\), the lower bound uses
\(k=0\) and needs no additional envelope condition.

\subparagraph{Theorem F.9. Finite-baseline minimax
limit}\label{theorem-f.9.-finite-baseline-minimax-limit}

Under the model and initialization above, (F.49), and \[
L/n\to\ell\in(0,\infty),\qquad M/n\to\rho\in(0,\infty),
\] the minimax risk (F.13) satisfies \[
\boxed{
\mathfrak R_{n,L,M}\longrightarrow
\sup_{P\in\mathcal P_\delta}
V_{\mathrm{FB}}\!\left(
K;\frac{c_P}{\sqrt\ell},\frac{c_P}{\sqrt\rho}
\right).
}
\tag{F.50}
\]

If instead \(L/n\to\infty\), with \(M/n\to\rho\in(0,\infty)\), then \[
\boxed{
\mathfrak R_{n,L,M}\longrightarrow
\sup_{P\in\mathcal P_\delta}V_G(K,c_P/\sqrt\rho).
}
\tag{F.51}
\]

This is an equality of prediction minimax values. It is not an assertion
that the full infinite-dimensional experiment is equivalent to a scalar
Gaussian experiment.

\textbf{Proof: upper bound for finite \(\ell\).} The constants \(c_P\)
lie in a fixed positive compact interval \(\mathcal C\). For every
\(\epsilon>0\) and every \(c\in\mathcal C\), choose an
\(\epsilon\)-optimal Gaussian rule for
\(V_{\mathrm{FB}}(K;c/\sqrt\ell,c/\sqrt\rho)\).

Such rules can be taken bounded Lipschitz with arbitrarily small extra
risk. Indeed, replace a measurable rule \(d\) by \(d_h(x)=E_Wd(x+hW)\).
Jensen and the TV distance between \(N(S,b^2)\) and \(N(S+hW,b^2)\),
conditional on \(S,W\), give \[
R_t(d_h)\le R_t(d)+\frac{h}{\pi b}.
\tag{F.52}
\] Gaussian smoothing also gives a finite Lipschitz constant.

The joint Gaussian law of \((S,D)\) depends continuously in TV on \(c\),
uniformly on \(\mathcal C\) and in \(t\). By Proposition F.7 the value
function is continuous as well. Choose finitely many nearby \(c\)-values
covering \(\mathcal C\), smooth a near-optimal rule at each, and
interpolate the rules using a Lipschitz partition of unity in \(c\).
Convexity of the loss gives a bounded jointly Lipschitz rule
\(d_\epsilon(c,d)\) with risk at most \[
V_{\mathrm{FB}}(K;c/\sqrt\ell,c/\sqrt\rho)+C\epsilon
\] at every \(c\in\mathcal C\).

Use the feasible predictor \[
\widehat\pi_\epsilon
=d_\epsilon(\widehat c,\widehat D_n).
\tag{F.53}
\] No independence between the scale estimate and the displacement
estimate is needed. For fixed \(\epsilon\) the interpolated rule has a
finite common Lipschitz constant \(J_\epsilon\), so replacing its
estimated first argument by c\_P changes expected loss by at most
\(J_\epsilon E|\widehat c-c_P|\), uniformly in P,k. Here \(\widehat c\)
denotes the truncated scale estimate. Only after taking n to infinity
for this fixed rule do we send epsilon to zero; we do not assert a
uniform Lipschitz bound for all near-optimal rules. Uniform consistency
of \(\widehat c\), the joint convergence in Theorem F.4 and the
conditional exit approximation imply \[
\limsup_n\mathfrak R_{n,L,M}
\le
\sup_PV_{\mathrm{FB}}(K;c_P/\sqrt\ell,c_P/\sqrt\rho)
+C\epsilon.
\] Let \(\epsilon\downarrow0\). A sufficiently slow diagonal sequence
gives asymptotically minimax feasible rules. In finite implementation,
the Gaussian rule is selected using the known design ratios \(L/n,M/n\);
no true baseline constants are supplied.

\textbf{Proof: lower bound for finite \(\ell\).} Fix any interior
baseline \(P_0\), any finite location window \(H\), and \(0\le K'<K\).
The subfamily (F.41) is eventually contained in the full experiment. To
check (F.49) at an endpoint, the strict inequality \(K'<K\) supplies the
envelope slack needed for (F.42); its actual seed is \(t+O(n^{-1/2})\).
Lemma F.8 yields \[
\liminf_n\mathfrak R_{n,L,M}
\ge V_H(K';c_0/\sqrt\ell,c_0/\sqrt\rho).
\] Proposition F.6 sends \(H\to\infty\). Continuity in \(K\) then sends
\(K'\uparrow K\). If \(K=0\), use the subfamily \(t=0\) directly, so the
same conclusion holds.

Interior baseline densities are dense in \(\mathcal P_\delta\) for the
required constants: replace any multiplier \(b\) by \((1-\zeta)b\), then
let \(\zeta\downarrow0\). Uniform monotonicity of the root equation and
bounded moments give continuity of \(q,\lambda,s,c\) along this path.
Thus taking the supremum over interior \(P_0\) gives the full right side
of (F.50), completing the lower bound.

\textbf{Proof for \(\ell=\infty\).} Theorem F.4 now has \(S=t\), so the
same upper construction applies to the known-baseline Gaussian value,
with \(\widehat c\) still estimated from data. For the lower bound,
restrict the baseline to one fixed \(P_0\). Its sample is independent of
the target parameter and the initialization error is \(o_P(n^{-1/2})\).
The aligned target subfamily therefore has the original Gaussian
calibration experiment and deterministic limiting target \(\Phi(t)\).
The same bounded-window likelihood argument, or its one-sample
specialization in Lemma F.8, yields the lower value
\(V_G(K,c_0/\sqrt\rho)\). Supremizing over \(P_0\) proves (F.51).
\(\square\)

\subparagraph{Corollary F.10. The information
threshold}\label{corollary-f.10.-the-information-threshold}

Under (F.49), suppose \(L/n\to\ell\in(0,\infty]\).

\begin{itemize}
\tightlist
\item
  If \(K>0\), uniformly consistent conditional prediction is possible if
  and only if \(M/n\to\infty\).
\item
  If \(K=0\) and \(\ell<\infty\), the same threshold holds over the
  unknown-baseline class.
\item
  If \(K=0\) and \(\ell=\infty\), the constant forecast \(1/2\) is
  uniformly consistent, without any target calibration.
\end{itemize}

\textbf{Proof.} Sufficiency in the first two cases follows from Theorem
F.4. For necessity, if \(M/n\) does not diverge, select a subsequence
with \(M/n\le C\). Give the predictor additional independent target
observations until \(M'=\lceil(C+1)n\rceil\). This can only decrease
minimax risk and does not change the target or initialization.

When \(\ell<\infty\), Theorem F.9 and the strictly positive lower bound
(F.35) give a positive limiting risk even at \(K=0\). When
\(\ell=\infty,K>0\), (F.51) gives a positive value: two distinct
\(t\)-values in the window have separated targets and overlapping
Gaussian observation densities. Thus the original risk cannot converge
to zero on this subsequence.

If \(K=0,\ell=\infty\), (F.23) gives \(S_n=o_P(1)\) uniformly, and
(F.21) gives the constant forecast assertion. \(\square\)

\subparagraph{Corollary F.11. A closed-form optimal rule for
root-preserving
shifts}\label{corollary-f.11.-a-closed-form-optimal-rule-for-root-preserving-shifts}

Suppose \(K=0\), \(L/n\to\ell\in(0,\infty)\), and
\(M/n\to\rho\in(0,\infty)\). Then the fully feasible rule

\[
\boxed{
\widehat\pi_{\mathrm{pool}}
=\Phi\!\left(\frac{M}{L+M}\widehat D_n\right)
}
\tag{F.54}
\]

is asymptotically minimax, with value

\[
\sup_{P\in\mathcal P_\delta}
\frac1\pi\arctan
\left[
\frac{(c_P/\sqrt\ell)(c_P/\sqrt\rho)}
{\sqrt{c_P^2/\ell+c_P^2/\rho+c_P^4/\ell^2}}
\right].
\tag{F.55}
\]

The weight \(M/(L+M)\) does not depend on the unknown dynamical scale.

\textbf{Proof.} In Proposition F.7, \[
\alpha=\frac{a_P^2}{a_P^2+b_P^2}
=\frac{\rho}{\ell+\rho}.
\] The map \((\alpha,d)\mapsto\Phi(\alpha d)\) has a bounded derivative
in \(\alpha\) on the positive compact range of design weights:
\(|d|\phi(\alpha d)\) is uniformly bounded there. Thus the finite design
weight \(M/(L+M)\), which converges to \(\alpha\), transfers the same
Gaussian risk uniformly without requiring a rate for convergence of L/n
or M/n.~Theorem F.4 transfers this risk uniformly, and Theorem F.9
supplies the matching lower bound.

There is also a direct pooling interpretation. When \(t_P(k)=0\), the
target satisfies \(v_Q(q_P)=0\) exactly under (F.4), so baseline and
target have the same selected root, even if \(k\ne0\). Let
\(\widehat q_{\mathrm{pool}}\) solve the empirical root equation using
all \(L+M\) observations, with the same bracket and fallback convention.
Since the baseline empirical drift is zero at \(x_L\) on the good event,
Taylor expansion gives \[
\frac{\sqrt n(x_L-\widehat q_{\mathrm{pool}})}{\widehat s}
=\frac{M}{L+M}\widehat D_n+o_P(1)
\] uniformly. Thus the pooled-root forecast is asymptotically the same
rule. Importantly, the future SGD still starts at \(x_L\); the pooled
root is used only in the forecast. \(\square\)

\paragraph{8. What this extension does and does not
establish}\label{what-this-extension-does-and-does-not-establish}

\subparagraph{8.1. A genuine unknown baseline, with a controlled
reference
class}\label{a-genuine-unknown-baseline-with-a-controlled-reference-class}

The unknown \(P\) ranges over a fixed infinite-dimensional density
neighborhood; it is not merely a Gaussian variance parameter. Its root
can differ from the known reference root by a fixed small amount. The
forecast estimates its root-dependent slope and noise scale from
\(B_L\).

The reference neighborhood, common risk and root bracket remain known.
Thus the result removes exact baseline knowledge within a declared
regular class. It does not remove every oracle component of the original
risk-isolation experiment.

The Gaussian reference supplies a concrete legal construction. The
proofs use uniform drift smoothness, a positive root derivative, nonzero
root noise, a common safe-region geometry, and a density lower bound for
finite exit. They do not assume the actual unknown baseline has Gaussian
class conditionals.

\subparagraph{8.2. Conditional target versus averaging over
initialization}\label{conditional-target-versus-averaging-over-initialization}

Equation (F.21) also implies, for each fixed \(u>0\), \[
\sup_{P,k}\Pr_B\{
|\pi_n(P,k;B_L)-\Phi(S_n)|>u
\}\longrightarrow0.
\] It does not assert a bound for every possible realized baseline
sample.

At \(K=0\), the marginal target approaches \(1/2\), but the conditional
target is approximately \(\Phi(a_PZ_L)\). The positive value in (F.35)
measures uncertainty about this realized conditional probability when
the baseline root is unknown. Averaging the initialization before
defining the loss would remove precisely this difficulty.

\subparagraph{8.3. Budget interpretation}\label{budget-interpretation}

For \(L\asymp n\), the baseline error is of the same order as the
stochastic seed of future SGD. This does not force \(L/n\to\infty\),
since that error is part of the realized target and can be measured
through target drift at the realized initialization.

Finite \(L\) still changes the optimal decision rule and its minimax
value. At \(K=0\), baseline and target observations estimate a common
root, giving the explicit pooling rule. For \(K>0\), pooling without
accounting for the bounded unknown root displacement can be biased; the
posterior-median problem (F.33) describes the appropriate tradeoff.

When \(L/n\to0\), the compact-seed argument in this chapter no longer
applies. The Gaussian model alone suggests that probabilities may
saturate and prediction may become easier. For example, its raw upper
bound tends to zero when \(a\to\infty\) and \(b/a\to0\). This is not an
actual-SGD theorem for that regime, and Corollary F.10 must not be
extrapolated to it.

\subparagraph{8.4. What still requires numerical or experimental
work}\label{what-still-requires-numerical-or-experimental-work}

The general-\(K\) value is characterized by (F.29)/(F.33), with matching
asymptotic lower and upper bounds. This chapter does not compute a
least-favorable prior or produce new numerical continuous-window
certificates for it. Proposition F.7 supplies a direct starting point
for that task.

The scale estimates are feasible and the raw and \(K=0\) optimal rules
are explicit. No actual-SGD confirmation is claimed here. Such
confirmation must resample the baseline sample as an outer level, hold
each realized initialization fixed across predictor comparisons, and
retain reference-trajectory uncertainty conditional on that
initialization.

The original optimized fixed-window rule is not automatically valid in
this new experiment. Equation (F.40) is a structural warning against
transferring its restricted output range.

\paragraph{9. Gaussian checks}\label{gaussian-checks}

These checks are diagnostic evaluations of exact formulas, not
statistical certificates for actual SGD.

{\def\LTcaptype{table} % do not increment counter
\begin{longtable}[]{@{}
  >{\raggedright\arraybackslash}p{(\linewidth - 6\tabcolsep) * \real{0.2500}}
  >{\raggedright\arraybackslash}p{(\linewidth - 6\tabcolsep) * \real{0.2500}}
  >{\raggedright\arraybackslash}p{(\linewidth - 6\tabcolsep) * \real{0.2500}}
  >{\raggedright\arraybackslash}p{(\linewidth - 6\tabcolsep) * \real{0.2500}}@{}}
\toprule\noalign{}
\begin{minipage}[b]{\linewidth}\raggedright
Initialization SD \(a\)
\end{minipage} & \begin{minipage}[b]{\linewidth}\raggedright
Calibration SD \(b\)
\end{minipage} & \begin{minipage}[b]{\linewidth}\raggedright
Exact \(K=0\) optimal MAE
\end{minipage} & \begin{minipage}[b]{\linewidth}\raggedright
Exact raw-rule maximum MAE
\end{minipage} \\
\midrule\noalign{}
\endhead
\bottomrule\noalign{}
\endlastfoot
0.5 & 1 & 0.1309898804 & 0.2322795272 \\
1 & 1 & \(1/6=0.1666666667\) & 0.1959132760 \\
2 & 0.5 & 0.0696044873 & 0.0700243480 \\
\end{longtable}
}

An independent numerical check used 1,500,000 Gaussian pairs and seed
20261030. The empirical pooled-rule MAEs were respectively 0.13108073,
0.16672833 and 0.06955975, with Monte Carlo standard errors 0.00007564,
0.00010471 and 0.00006992. All agree with their analytic values within
two Monte Carlo standard errors. These errors describe the numerical
check only.

A minimal reproduction is:

\begin{lstlisting}[language=Python]
import numpy as np
from scipy.special import ndtr
from math import atan, pi, sqrt

rng = np.random.default_rng(20261030)
n_draws = 1_500_000
z = rng.standard_normal((2, n_draws))

for a, b in [(0.5, 1.0), (1.0, 1.0), (2.0, 0.5)]:
    seed = a * z[0]
    observation = seed + b * z[1]
    alpha = a * a / (a * a + b * b)
    loss = np.abs(ndtr(alpha * observation) - ndtr(seed))
    exact = atan(a * b / sqrt(a*a + b*b + a**4)) / pi
    print(a, b, exact, loss.mean(), loss.std() / sqrt(n_draws))
\end{lstlisting}

\subsubsection{Experimental interface}\label{experimental-interface}

Actual-SGD confirmation must resample baseline observations at an outer
level, hold each realized initialization fixed across predictor
comparisons, and preserve conditional reference-trajectory uncertainty.
The pooling rule is optimal only for the declared scalar root-preserving
subcase. Finite-horizon BN risk, batch-statistics coupling and omitted
directions require separate validation.

\subsubsection{Classical background and proof
attribution}\label{classical-background-and-proof-attribution}

\begin{itemize}
\tightlist
\item
  Project source: the supplied manuscript, \emph{How Much Unlabeled
  Information Is Needed to Predict Entropy-SGD Risk Exit?}, especially
  its actual-SGD proof in Section 4, likelihood comparison in Section 5,
  and bounded unknown-direction argument in Section 6. This chapter
  repeats the relevant stopped-seed estimates and gives the additional
  initialization and marked-loss arguments explicitly.
\item
  Yuri Bakhtin, \emph{Exit asymptotics for small diffusion about an
  unstable equilibrium}, Stochastic Processes and their Applications
  118(5), 839--851, 2008.
  \href{https://arxiv.org/abs/math/0701569}{Author preprint}. Background
  for the classical unstable-exit mechanism; its diffusion theorem is
  not substituted for the fresh-batch discrete proof here.
\item
  A. W. van der Vaart, \emph{Asymptotic Statistics}, Cambridge
  University Press, 1998, especially Chapter 9, \emph{Limits of
  Experiments}.
  \href{https://www.cambridge.org/core/books/abs/asymptotic-statistics/limits-of-experiments/3DCF37272E376213D53F0CABB3EB9CC5}{Publisher
  chapter}. Background for experiment comparison; Lemma F.8 supplies the
  initialization-preserving argument needed for this sample-dependent
  loss.
\end{itemize}

No exhaustive novelty or priority search is claimed. The external
sources above provide background attribution, not proofs of the new
model-specific statements.
